% 8_evaluation.tex — 模型评价

\subsection{模型优点}

\begin{enumerate}[label=(\arabic*)]
    \item \textbf{几何模型严谨}：基于视线遮挡的解析几何判定具有严格的数学基础，通过目标圆柱离散化实现了精度与效率的平衡，避免了经验性遮挡判定的主观性。

    \item \textbf{分层优化策略合理}：针对不同复杂度的子问题采用递进的求解策略（直接计算$\to$网格搜索$\to$遗传算法），充分利用了各算法的优势，兼顾了全局搜索能力和局部收敛精度。

    \item \textbf{模块化程度高}：MATLAB代码采用运动学$\to$几何$\to$优化的三层模块化架构，各模块接口清晰、可独立测试，便于扩展到更多无人机和导弹的场景。

    \item \textbf{物理建模完整}：考虑了导弹的直线飞行、无人机的等高度匀速飞行、干扰弹的自由落体和烟幕云团的下沉运动，物理过程描述完整。

    \item \textbf{灵敏度分析充分}：对烟幕半径、无人机速度、导弹速度、下沉速度等关键参数进行了系统的灵敏度分析，揭示了各参数对遮蔽效果的影响规律。

    \item \textbf{多导弹均衡优化}：问题5采用maximin准则，确保对各导弹均有较好的遮蔽效果，避免了过度优化单一导弹而忽视其他威胁的问题。
\end{enumerate}

\subsection{模型缺点}

\begin{enumerate}[label=(\arabic*)]
    \item \textbf{烟幕模型简化}：实际烟幕云团的扩散受大气湍流、风速风向等气象条件影响，并非严格的球形和匀速下沉。本文的确定性球状模型未考虑这些随机因素。

    \item \textbf{遮蔽判据二值化}：模型假设"完全遮蔽"或"完全不遮蔽"的二值判断，而实际中烟幕浓度从中心向外递减，遮蔽效果存在渐变过渡区。

    \item \textbf{忽略无人机转弯过程}：假设无人机可瞬时调整航向，而实际中无人机有转弯半径和加减速过程的限制。

    \item \textbf{未考虑导弹机动}：假设导弹始终飞向假目标且不进行末端机动，而实际中导弹可能具有目标识别和规避机动的能力。

    \item \textbf{计算资源限制}：问题5的40维优化在有限计算资源下难以保证达到全局最优，实际求解依赖启发式分解策略。
\end{enumerate}

\subsection{改进方向}

\begin{enumerate}[label=(\arabic*)]
    \item \textbf{引入随机烟幕扩散模型}：采用高斯烟团扩散模型或计算流体力学(CFD)方法描述烟幕在大气中的实际扩散过程，提高仿真逼真度\cite{zhang2020}。

    \item \textbf{连续遮蔽度模型}：将二值遮挡判定扩展为基于烟幕浓度分布的连续遮蔽度函数，更准确地反映遮蔽效果的空间分布。

    \item \textbf{动态重规划能力}：引入模型预测控制(MPC)或强化学习(RL)方法，使无人机能够根据实时态势动态调整飞行策略，应对导弹机动等不确定因素。

    \item \textbf{多目标优化框架}：将问题扩展为多目标优化问题，同时考虑遮蔽时间、无人机燃油消耗、干扰弹使用数量等多个优化目标，提供帕累托最优解集供决策者选择\cite{gao2009}。

    \item \textbf{并行计算加速}：利用MATLAB的并行计算工具箱对GA的适应度评估进行并行化，提高高维问题的求解效率。
\end{enumerate}
